\itemhead{2.1 例题 · 例题6（教材p.39）}{XJ-B1-C02-EX-P039-Y0870}
\label{XJ-B1-C02-EX-P039-Y0870}
{\footnotesize 湘教版教材p.39 / PDF第46页 / 2.1 例题 / 例题6}
\field{原题}{对任意三个正实数 $a,b,c$，求证：$a+b+c\ge\sqrt{ab}+\sqrt{bc}+\sqrt{ca}$，当且仅当 $a=b=c$ 时等号成立。}
\field{参考答案}{分别写 $a+b\ge2\sqrt{ab}$、$b+c\ge2\sqrt{bc}$、$c+a\ge2\sqrt{ca}$，相加除以2即可；等号当且仅当 $a=b=c$。}
{\small \field{关联知识}{I03 基本不等式、等号与最值建模（教材帮p.52、53、54 / PDF61、62、63）}
\field{方法与易错点}{列正数条件与固定量，构造基本不等式，解等号条件并代回实际限制。 易错：把一个下界误当最小值；多个等号条件不能同时实现；漏单位和整数限制。}
\field{答案依据}{例题原解位置：湘教版教材p.39及题库原图；独立编写可编辑参考解答。答案册此部分主要收练习、习题，不伪造例题答案册位置。}}
